% part: intuitionistic-logic
% chapter: soundness-completeness
% section: lindenbaum

\documentclass[../../../include/open-logic-section]{subfiles}

\begin{document}

\olfileid{int}{sc}{lin}

\olsection{লিন্ডেনবাউমের উপপাদ্য}

স্বজ্ঞাবাদী যুক্তিবিদ্যার সম্পূর্ণতার উপপাদ্য প্রমাণে $\Gamma \Proves/ !A$
ধরে এমন একটি মডেল নির্মাণ করি যাতে $\mSat{M}{\Gamma}$ অথচ
$\mSat/{M}{!A}$।

ধ্রুপদি যুক্তিবিদ্যায় !!{derivability}-র সম্পর্কটিকে সঙ্গতির ধারণায়
নামিয়ে আনা যায়: !!{formula}~$!A$ একটি !!{formula}s-এর সেট থেকে
!!{derivable} হবে তখনই, যখন সেটটির সঙ্গে~$!A$-এর নঞর্থক রূপ
যোগ করলে সেটটি অসঙ্গত হয়। স্বজ্ঞাবাদী যুক্তিবিদ্যায় তা সম্ভব নয়।
এখানে $\lnot!A$ অসঙ্গত হলে কেবল $\Proves \lnot\lnot !A$ পাওয়া যায়।
সাধারণভাবে $\lnot\lnot!A \lif !A$ স্বজ্ঞাবাদীভাবে সিদ্ধ নয়; তাই
$\Proves !A$ সিদ্ধান্তে পৌঁছানো যায় না।

সুতরাং মডেল~$\mModel{M}$ নির্মাণের সময় !!{formula}~$!A$-এর
অ-!!{derivability}-র হিসাব রাখতে হবে। ফলে ধ্রুপদি প্রমাণের মতো
সম্পূর্ণ সেট $\Gamma^* \supseteq \Gamma$ দিয়ে মডেল~$\mModel{M}$
বানানো যাবে না: প্রতিটি সম্পূর্ণ সেট~$\Gamma^*$-তে
$\Gamma^* \Proves !A \lor \lnot !A$।

সম্পূর্ণ সেট~$\Gamma^*$-এর বদলে আমরা সূত্রসমূহের একটি
\emph{প্রাইম সেট} ব্যবহার করব:

\begin{defn}\ollabel{defn:prime}
  !!{formula}s-এর সেট~$\Gamma$ \emph{প্রাইম} তখনই, যখন
  \begin{enumerate}
  \item\ollabel{defn:prime1} $\Gamma$ সঙ্গত, অর্থাৎ $\Gamma \Proves/ \lfalse$;
  \item\ollabel{defn:prime2} $\Gamma \Proves !A$ হলে $!A \in \Gamma$; এবং
  \item\ollabel{defn:prime3} $!A \lor !B \in \Gamma$ হলে $!A \in
    \Gamma$ অথবা $!B \in \Gamma$।
  \end{enumerate}
\end{defn}

\begin{lem}[লিন্ডেনবাউমের উপপাদ্য]
  \ollabel{lem:lindenbaum} $\Gamma \Proves/ !A$ হলে এমন একটি
  $\Gamma^* \supseteq \Gamma$ আছে, যাতে $\Gamma^*$ প্রাইম এবং
  $\Gamma^* \Proves/ !A$।
\end{lem}

\begin{proof}
  $!B_1 \lor !C_1$, $!B_2 \lor !C_2$, \dots, দিয়ে
  $!B \lor !C$-রূপের সমস্ত !!{formula}s-এর একটি গণনা ধরি।
  !!{formula}s-এর সেটের একটি বর্ধমান ক্রম~$\Gamma_n$ সংজ্ঞায়িত করব:
  প্রতিটি $\Gamma_{n+1}$ পাওয়া যাবে $\Gamma_n$-এর সঙ্গে একটি নতুন
  !!a{formula} যোগ করে। সব~$\Gamma_n$-এর সংযোগই হবে $\Gamma^*$।
  নতুন !!{formula}s এমনভাবে বাছব যাতে $\Gamma^*$ প্রাইম হয়, অথচ
  $\Gamma^* \Proves/ !A$ বজায় থাকে। প্রতি ধাপে তাই প্রথম এমন
  বিযুক্তি~$!B_i \lor !C_i$ খুঁজব যার জন্য:
  \begin{enumerate}
  \item\ollabel{gamma-1} $\Gamma_n \Proves !B_i \lor !C_i$
  \item\ollabel{gamma-2} $!B_i \notin \Gamma_n$ এবং $!C_i \notin \Gamma_n$
  \end{enumerate}
  $\Gamma_n \cup \{!B_i\} \Proves/ !A$ হলে $!B_i$, অন্যথায়
  $!C_i$ যোগ করব। এই পদ্ধতি যে কাজ করে, তা দেখাতে হবে। আপাতত
  \olref{gamma-1} ও~\olref{gamma-2} পূরণকারী ক্ষুদ্রতম $i$-কে
  $i(n)$ বলি।
  $i(n)$-এর মান $i$ নির্ধারিত হলে $\Gamma_n$-র প্রথম অনিষ্পন্ন
  বিযুক্তি $!B_i \lor !C_i$; তার $!B_i$ বা $!C_i$ যোগ হয়।

  $\Gamma_0 = \Gamma$ ধরি এবং
  \[
  \Gamma_{n+1} = \begin{cases}
    \Gamma_n \cup \{!B_{i(n)}\} &
    \text{if $\Gamma_n \cup \{!B_{i(n)}\} \Proves/ !A$} \\
    \Gamma_n \cup \{!C_{i(n)}\} & \text{otherwise}
    \end{cases}
  \]
  যদি $i(n)$ অনির্দিষ্ট হয়—অর্থাৎ $\Gamma_n \Proves !B \lor !C$
  হলেই $!B \in \Gamma_n$ অথবা $!C \in \Gamma_n$—তবে
  $\Gamma_{n+1} = \Gamma_n$ ধরি। এবার
  $\Gamma^* = \bigcup_{n=0}^\infty \Gamma_n$ নিই।

  প্রথমে দেখাই, সব $n$-এর জন্য $\Gamma_n \Proves/ !A$।
  $n$-এর ওপর আরোহে এগোই। $n = 0$-তে দাবিটি উপপাদ্যের অনুমান
  $\Gamma \Proves/ !A$ থেকেই আসে। $n>0$ হলে দেখাতে হবে:
  $\Gamma_n \Proves/ !A$ থেকে $\Gamma_{n+1} \Proves/ !A$।
  $i(n)$ অনির্দিষ্ট হলে $\Gamma_{n+1} = \Gamma_n$, তাই আর কিছু
  দেখাতে হয় না। অন্যথায়, সংক্ষেপে $i = i(n)$ ধরি।

  দাবিটির বিপরীতার্থক রূপ প্রমাণ করি। ধরি $\Gamma_{n+1}
  \Proves !A$। নির্মাণ অনুযায়ী, $\Gamma_n \cup
  \{!B_i\} \Proves/ !A$ হলে $\Gamma_{n+1} = \Gamma_n \cup
  \{!B_i\}$; অন্যথায় $\Gamma_{n+1} = \Gamma_n \cup \{!C_i\}$।
  প্রথমটি হতে পারে না, কারণ সেক্ষেত্রে $\Gamma_{n+1} \Proves/ !A$।
  অতএব $\Gamma_n \cup \{!B_i\} \Proves !A$ এবং
  $\Gamma_{n+1} = \Gamma_n \cup \{!C_i\}$। $i(n)$-এর সংজ্ঞায়
  $\Gamma _n \Proves !B_i \lor !C_i$। আবার
  $\Gamma_n \cup \{!B_i\} \Proves !A$ এবং
  $\Gamma_{n+1} = \Gamma_n \cup \{!C_i\} \Proves !A$।
  তাই $\Gamma_n \Proves !A$; এটাই চেয়েছিলাম।

  $\Gamma^* \Proves !A$ হলে কোনো সসীম উপসেট $\Gamma'
  \subseteq \Gamma^*$ থেকে $\Gamma' \Proves !A$ হতো। প্রতিটি
  $!D \in \Gamma'$ কোনো-না-কোনো~$\Gamma_i$-তে থাকে। সেই সূচকগুলির
  বৃহত্তমটি $n$ ধরি। $i \le n$ হলে $\Gamma_i \subseteq \Gamma_n$;
  সুতরাং $\Gamma' \subseteq \Gamma_n$। কিন্তু এতে
  $\Gamma_n \Proves !A$ হয়, যা আগের ফল
  $\Gamma_n \Proves/ !A$-এর বিরোধী।

  শেষে দেখাই, $\Gamma^*$ প্রাইম—অর্থাৎ
  \olref{defn:prime}-এর \olref{defn:prime1},
  \olref{defn:prime2} এবং~\olref{defn:prime3} শর্তগুলি পূরণ করে।

  প্রথমত, $\Gamma^* \Proves/ !A$; তাই $\Gamma^*$ সঙ্গত,
  অর্থাৎ \olref{defn:prime1} পূরণ করে।

  এবার দেখাই, $\Gamma^* \Proves !B \lor !C$ হলে
  $!B \in \Gamma^*$ অথবা $!C \in \Gamma^*$। তাতেই
  \olref{defn:prime3} প্রতিষ্ঠিত হবে: $!B \lor !C \in \Gamma^*$
  হলে অবশ্যই $\Gamma^* \Proves !B \lor !C$। তাই ধরি,
  $\Gamma^* \Proves !B \lor !C$, অথচ $!B \notin \Gamma^*$ এবং
  $!C \notin \Gamma^*$। কোনো~$n$-এর জন্য
  $\Gamma_n \Proves !B \lor !C$। সমস্ত বিযুক্তির গণনায়
  $!B \lor !C$-টি $!B_j \lor !C_j$ হিসেবে আসে। তাই
  $!B_j \lor !C_j$ সূচক~$i(n)$-এর সংজ্ঞার শর্ত পূরণ করে:
  $\Gamma_n \Proves !B_j \lor !C_j$, কিন্তু
  $!B_j \notin \Gamma_n$ ও $!C_j \notin \Gamma_n$।
  % BN-SRC-405: only the finitely many eligible indices below fixed j must be exhausted.
  ওই সূচকের চেয়ে ছোট সূচক সসীমসংখ্যক। এদের কোনোটি একবার বেছে নিয়ে
  তার একটি বিযোজ্য পদ যোগ করলে সেই সূচক আর কখনো দ্বিতীয় শর্ত
  পূরণ করতে পারে না। অথচ নির্দিষ্ট সূচকটির জন্য শর্ত দুটি পূরণ হতে থাকে:
  সূত্রনিষ্পাদন বর্ধমান থাকে এবং কোনো বিযোজ্য পদই যুক্ত হয় না।
  ফলে কোনো ধাপ~$m$-এ $j = i(m)$ হতেই হবে। তখন হয়
  $!B \in \Gamma_{m+1}$, নয় $!C \in \Gamma_{m+1}$;
  এটি $!B \notin \Gamma^*$ ও $!C \notin \Gamma^*$-এর বিরোধী।
  তাই $\Gamma^* \Proves !B \lor !C$ অনুমানটির দাবি প্রতিষ্ঠিত।

  এখন ধরি $\Gamma^* \Proves !B$। তাহলে
  $\Gamma^* \Proves !B \lor !B$। এইমাত্র প্রমাণ করেছি,
  $\Gamma^* \Proves !B \lor !B$ হলে $!B \in \Gamma^*$।
  সুতরাং $\Gamma^*$, \olref{defn:prime}-এর
  \olref{defn:prime2} শর্তও পূরণ করে।
\end{proof}

\begin{prob}
  দেখান, $\Gamma \Proves/ \bot$ হলে ধ্রুপদি যুক্তিবিদ্যায়
  $\Gamma$ সঙ্গত; অর্থাৎ এমন একটি মানায়ন আছে যাতে
  $\Gamma$-র সব সূত্র সত্য।
\end{prob}

\end{document}
